\section{FOC for \texorpdfstring{$\underline{r}^{*}$}{underline r star}}
\label{app:foc-low}

This appendix derives \eqref{eq:FOC-low-simplified} from \eqref{eq:Pi-low-total}.
Write $z\equiv\beta^{-1}(\underline{r})$ and $v_{\psi}$ from \eqref{eq:v-psi-def}.
Total profit splits into the failure and success components
\eqref{eq:profit-fail-closed} and \eqref{eq:profit-success}:
\begin{equation}
  \Pi(\underline{r})
  = \mathbf{1}\{\underline{r}>\underline{r}_0\}\,
    \Pi_{\text{fail}}(\underline{r})
    + \Pi_{\text{succ}}(\underline{r}),
  \label{eq:Pi-low-split}
\end{equation}
where $\Pi_{\text{fail}}=0$ for $\underline{r}\le\underline{r}_0$.
Differentiating \eqref{eq:Pi-low-split} (holding fixed the region
$\underline{r}>\underline{r}_0$ where the failure formula applies) gives
\begin{equation}
  \frac{\mathrm d\Pi(\underline{r})}{\mathrm d\underline{r}}
  = \mathbf{1}\{\underline{r}>\underline{r}_0\}\,
    \frac{\mathrm d\Pi_{\text{fail}}}{\mathrm d\underline{r}}
    + \frac{\mathrm d\Pi_{\text{succ}}}{\mathrm d\underline{r}}.
  \label{eq:dPi-total}
\end{equation}

\subsection{Derivative of \texorpdfstring{$\Pi_{\text{fail}}$}{Pi fail}}
\label{subsec:deriv-fail}

For $\underline{r}>\underline{r}_0$, $\Pi_{\text{fail}}$ is
\eqref{eq:profit-fail-closed}.
Since $z=(\underline{r}-\delta\psi c_L)/(1-\delta\psi)$,
$\mathrm dz/\mathrm d\underline{r}=1/(1-\delta\psi)$.
Differentiating,
\begin{align}
  \frac{\mathrm d\Pi_{\text{fail}}}{\mathrm d\underline{r}}
  &= \frac{d}{d\underline{r}}
    \left[
      \delta N (c_H - c_L)\,
      \bigl(F(z) - F(v_{\psi})\bigr)\,
      F(z)^{N-1}
    \right]
  \notag\\
  &= \delta N (c_H - c_L)\,
    \frac{1}{1 - \delta \psi}\,
    f(z)\, F(z)^{N-2}
    \Bigl[
      N F(z)
      - (N-1) F(v_{\psi})
    \Bigr].
  \label{eq:deriv-fail}
\end{align}

\subsection{Derivative of \texorpdfstring{$\Pi_{\text{succ}}$}{Pi succ}}
\label{subsec:deriv-succ}

Write $\Pi_{\text{succ}}$ as in \eqref{eq:profit-success}:
\begin{equation}
  \Pi_{\text{succ}}
  = N \Biggl\{
    (\underline{r} - c_L) F(\underline{r})^{N-1}
      \bigl(1 - F(z)\bigr)
    + \int_{z}^{1}
      \int_{\underline{r}}^{v} (s - c_L)\, dF(s)^{N-1}
      \, dF(v)
  \Biggr\}.
  \label{eq:Pi-succ-appendix}
\end{equation}
The outer bracket and the double integral are differentiated separately.

\paragraph{Double integral.}
Integration by parts gives
\begin{align}
  &\int_{z}^{1}
    \int_{\underline{r}}^{v} (s - c_L)\, dF(s)^{N-1}
    \, dF(v)
  \notag\\
  =&\; \int_{\underline{r}}^{1} (s - c_L)\, dF(s)^{N-1}
    - F(z)
      \int_{\underline{r}}^{z}
        (s - c_L)\, dF(s)^{N-1}
    - \int_{z}^{1}
      F(v)(v - c_L)\, dF(v)^{N-1}.
  \label{eq:double-int-parts}
\end{align}
Differentiating term by term yields \eqref{eq:deriv-int1}--\eqref{eq:deriv-int3}
and the collapsed form
\begin{equation}
  \frac{d}{d\underline{r}}
    \int_{\underline{r}}^{1} (s - c_L)\, dF(s)^{N-1}
  = -(N-1)(\underline{r} - c_L) F(\underline{r})^{N-2} f(\underline{r}).
  \label{eq:deriv-int1}
\end{equation}
\begin{align}
  &\frac{d}{d\underline{r}}
    \left(
      F(z)
      \int_{\underline{r}}^{z}
        (s - c_L)\, dF(s)^{N-1}
    \right)
  \notag\\
  =&\; f(z)\,
    \frac{\mathrm dz}{\mathrm d\underline{r}}
    \int_{\underline{r}}^{z}
      (s - c_L)\, dF(s)^{N-1}
  \notag\\
  &\quad + (N-1) F(z)
    \Bigl[
      (z - c_L) F(z)^{N-2} f(z)\,
        \frac{\mathrm dz}{\mathrm d\underline{r}}
      - (\underline{r} - c_L) F(\underline{r})^{N-2} f(\underline{r})
    \Bigr].
  \label{eq:deriv-int2}
\end{align}
\begin{equation}
  \frac{d}{d\underline{r}}
    \int_{z}^{1}
      F(v)(v - c_L)\, dF(v)^{N-1}
  = -\frac{\mathrm dz}{\mathrm d\underline{r}}
    (N-1)(z - c_L) F(z)^{N-1} f(z).
  \label{eq:deriv-int3}
\end{equation}
Combining \eqref{eq:deriv-int1}--\eqref{eq:deriv-int3} gives
\begin{align}
  &\frac{\mathrm d}{\mathrm d\underline{r}}
    \int_{z}^{1}
    \int_{\underline{r}}^{v} (s - c_L)\, dF(s)^{N-1}
    \, dF(v)
  \notag\\
  =&\; -(N-1)(\underline{r} - c_L) F(\underline{r})^{N-2} f(\underline{r})
    \bigl(1 - F(z)\bigr)
  \notag\\
  &\quad - f(z)\,
    \frac{d\beta^{-1}(\underline{r})}{d\underline{r}}
    \int_{\underline{r}}^{z}
      (s - c_L)\, dF(s)^{N-1}.
  \label{eq:deriv-double-int}
\end{align}

\paragraph{Bracket term.}
\begin{align}
  &\frac{d}{d\underline{r}}
    \left[
      (\underline{r} - c_L) F(\underline{r})^{N-1}
        \bigl(1 - F(z)\bigr)
    \right]
  \notag\\
  =&\; \Biggl[
    \begin{aligned}[t]
      &F(\underline{r})\bigl(1 - F(z)\bigr) \\
      &\quad + (N-1)(\underline{r} - c_L)
        \bigl(1 - F(z)\bigr) f(\underline{r}) \\
      &\quad - (\underline{r} - c_L) F(\underline{r})
        f(z)\,
        \frac{1}{1 - \delta \psi}
    \end{aligned}
  \Biggr]
  F(\underline{r})^{N-2}.
  \label{eq:deriv-bracket}
\end{align}

Multiplying \eqref{eq:deriv-bracket} and \eqref{eq:deriv-double-int} by $N$
gives $\mathrm d\Pi_{\text{succ}}/\mathrm d\underline{r}$.
The $(N-1)(\underline{r}-c_L)$ terms cancel, leaving
\begin{align}
  \frac{\mathrm d\Pi_{\text{succ}}}{\mathrm d\underline{r}}
  &= N F(\underline{r})^{N-2}
    \Bigl[
      F(\underline{r})\bigl(1 - F(z)\bigr)
      - (\underline{r} - c_L) F(\underline{r})\,
        f(z)\,
        \frac{\mathrm dz}{\mathrm d\underline{r}}
    \Bigr]
  \notag\\
  &\quad - N f(z)\,
    \frac{\mathrm dz}{\mathrm d\underline{r}}
    \int_{\underline{r}}^{z}
      (s - c_L)\, dF(s)^{N-1}.
  \label{eq:deriv-succ-collapsed}
\end{align}

\subsection{Full FOC and simplification}
\label{subsec:foc-full}

Substituting \eqref{eq:deriv-fail} and \eqref{eq:deriv-succ-collapsed} into
\eqref{eq:dPi-total} and setting $\mathrm d\Pi(\underline{r})/\mathrm d\underline{r}=0$
at $\underline{r}=\underline{r}^{*}$ gives
\begin{align}
  0
  &= \mathbf{1}\{\underline{r}^{*}>\underline{r}_0\}\,
    \delta N (c_H - c_L)\,
    \frac{1}{1 - \delta \psi}\,
    f(z^{*})\, F(z^{*})^{N-2}
    \Bigl[
      N F(z^{*})
      - (N-1) F(v_{\psi})
    \Bigr]
  \notag\\
  &\quad + N F(\underline{r}^{*})^{N-2}
    \Bigl[
      F(\underline{r}^{*})\bigl(1 - F(z^{*})\bigr)
      - (\underline{r}^{*} - c_L) F(\underline{r}^{*})\,
        f(z^{*})\,
        \frac{\mathrm dz}{\mathrm d\underline{r}}
    \Bigr]
  \notag\\
  &\quad - N f(z^{*})\,
    \frac{\mathrm dz}{\mathrm d\underline{r}}
    \int_{\underline{r}^{*}}^{z^{*}}
      (s - c_L)\, dF(s)^{N-1}.
  \label{eq:FOC-low-full}
\end{align}
Collecting the $(\underline{r}^{*}-c_L)$ terms on the left and dividing by
$N F(\underline{r}^{*})^{N-1} f(z^{*})\,\mathrm dz/\mathrm d\underline{r}$
yields \eqref{eq:FOC-low-simplified}.
Explicitly, \eqref{eq:FOC-low-full} is equivalent to
\begin{equation}
  N F(\underline{r}^{*})^{N-1} f(z^{*})\,
    \frac{\mathrm dz}{\mathrm d\underline{r}}\,
    (\underline{r}^{*} - c_L)
  = N F(\underline{r}^{*})^{N-1}\bigl(1 - F(z^{*})\bigr)
    - N f(z^{*})\,
      \frac{\mathrm dz}{\mathrm d\underline{r}}
      \int_{\underline{r}^{*}}^{z^{*}}
        (s - c_L)\, dF(s)^{N-1}
    + \mathbf{1}\{\underline{r}^{*}>\underline{r}_0\}\,
      \frac{\mathrm d\Pi_{\text{fail}}}{\mathrm d\underline{r}},
  \label{eq:FOC-collect-r}
\end{equation}
so
\begin{equation}
  \underline{r}^{*} - c_L
  = \frac{1 - F(z^{*})}{f(z^{*})\,\mathrm dz/\mathrm d\underline{r}}
    + \mathbf{1}\{\underline{r}^{*}>\underline{r}_0\}\,
      \frac{
        \mathrm d\Pi_{\text{fail}}/\mathrm d\underline{r}
      }{
        N F(\underline{r}^{*})^{N-1} f(z^{*})\,
          \mathrm dz/\mathrm d\underline{r}
      }
    - \frac{1}{F(\underline{r}^{*})^{N-1}}
      \int_{\underline{r}^{*}}^{z^{*}}
        (s - c_L)\, dF(s)^{N-1},
  \label{eq:FOC-unified}
\end{equation}
which is \eqref{eq:FOC-low-simplified}.
Substituting \eqref{eq:FOC-unified} back into \eqref{eq:FOC-low-full} shows
$\mathrm d\Pi(\underline{r}^{*})/\mathrm d\underline{r}=0$ if and only if
\eqref{eq:FOC-low-simplified} holds:
\begin{equation}
  \frac{\mathrm d\Pi}{\mathrm d\underline{r}}
  = -N F(\underline{r})^{N-1} f(z)\,
    \frac{\mathrm dz}{\mathrm d\underline{r}}\,
    \Bigl[
      (\underline{r} - c_L) - \text{RHS of \eqref{eq:FOC-low-simplified}}
    \Bigr].
  \label{eq:FOC-residual}
\end{equation}
The first line on the right-hand side of \eqref{eq:FOC-low-simplified} comes
from the success bracket in \eqref{eq:deriv-succ-collapsed} and is present
whether or not $\underline{r}^{*}>\underline{r}_0$.
The failure adjustment (second line) comes from \eqref{eq:deriv-fail}.
The integral is the success double-integral term.

\paragraph{Algebraic check.}
When $\mathbf{1}\{\underline{r}^{*}>\underline{r}_0\}=1$, all three lines
contribute.
When $\mathbf{1}\{\underline{r}^{*}>\underline{r}_0\}=0$,
\eqref{eq:FOC-low-simplified} reduces to the first-order condition from
$\mathrm d\Pi_{\text{succ}}/\mathrm d\underline{r}=0$ alone.
